\documentclass[11pt]{article}
\usepackage[margin=2.0cm]{geometry}
\usepackage{booktabs}
\usepackage{amsmath}
\usepackage{amssymb}
\usepackage{enumitem}
\usepackage{array}
\usepackage{listings}
\usepackage{xcolor}
\usepackage[T1]{fontenc}
\usepackage{lmodern}

\setlength{\parindent}{0pt}
\setlength{\parskip}{3.6pt}
\newcommand{\code}[1]{\texttt{#1}}
\newcommand{\hr}{\vspace{2pt}\hrule\vspace{6pt}}
\renewcommand{\arraystretch}{1.12}

\definecolor{kw}{RGB}{0,90,160}
\definecolor{cmt}{RGB}{110,110,110}
\lstdefinestyle{vhdl}{
  language=VHDL, basicstyle=\ttfamily\footnotesize,
  keywordstyle=\color{kw}\bfseries, commentstyle=\color{cmt}\itshape,
  showstringspaces=false, columns=fullflexible, breaklines=true,
  frame=leftline, framesep=6pt, xleftmargin=8pt, aboveskip=6pt, belowskip=6pt
}
\lstset{style=vhdl}

\begin{document}

\begin{center}
{\Large\bfseries EE\,362 --- Computer Organization Laboratory}\\[2pt]
{\large\bfseries Lab 5 --- Guided Examples: Solution}\\[2pt]
{\itshape Ex 1 one function unit \quad$\cdot$\quad Ex 2 what an ALU is \quad$\cdot$\quad Ex 3 the three shifts}\\[4pt]
Fall 2026 \quad$\cdot$\quad Lab TA: Gandhi Torres
\end{center}
\hr

\begin{center}
\begin{tabular}{clll}
\toprule
& circuit & new idea & feeds the assignment as \\
\midrule
1 & \code{alu\_add} & one function unit \textbf{plus flags} & the template all 16 operations copy \\
2 & \code{alu\_mini} & an ALU is \textbf{parallel units $+$ a mux} & why your \code{case} \emph{is} the mux \\
3 & \code{shift\_demo} & the \textbf{type} decides the fill & \code{SRL} vs \code{SRA} \\
\bottomrule
\end{tabular}
\end{center}

\vspace{4pt}\hr

\section*{Example 1 --- One function unit, complete}

Read this one carefully and the other fifteen operations become one line each.

\subsection*{1.1 The arithmetic, and why \code{unsigned} is used for both}

\begin{lstlisting}
result <= std_logic_vector( unsigned(Ain) + unsigned(Bin) );
\end{lstlisting}

\code{std\_logic\_vector} is a bag of bits with no numeric meaning, so you cannot
add two of them. Convert, add, convert back.

\begin{quote}
\textbf{We add as \code{unsigned} even when the values are signed, and that is
correct.} Two's complement addition is \textbf{bit-identical} whether you call the
operands signed or unsigned. Only the \emph{interpretation} of the result differs
--- and interpreting the result is what the flags are for.
\end{quote}

This is worth stating aloud because it is the one place where the type does
\emph{not} matter. In Example 3 the type is everything.

\subsection*{1.2 The flags}

\begin{align*}
\text{Zero} &= \overline{R_{31} + R_{30} + \cdots + R_0} && \text{a 32-input NOR: 1 only when every bit is 0}\\
\text{Neg}  &= R_{31} && \text{in two's complement the top bit \emph{is} the sign}
\end{align*}

\begin{lstlisting}
Zero <= '1' when result = (result'range => '0') else '0';
Neg  <= result(31);
\end{lstlisting}

\textbf{\code{Neg} is a wire, not a gate.} Nothing is computed. Students often
expect a comparison here; there is none to make.

Note \code{(result'range => '0')}: an all-zeros aggregate shaped like
\code{result}. It works at any width without being rewritten, and it is how you
compare a vector against zero.

\subsection*{1.3 Carry out, and why you must widen first}

Adding two 32-bit numbers can produce 33 bits. To see the 33rd, \textbf{widen the
operands before adding}:

\begin{lstlisting}
sum33 <= unsigned('0' & Ain) + unsigned('0' & Bin);   -- 33 bits wide
Cout  <= sum33(32);
\end{lstlisting}

Add at 32 bits and the extra bit is simply lost --- there is nowhere for it to go.
This is the same idea as Lab 4's extension: \textbf{make room first, then
compute.}

\vspace{4pt}\hr

\section*{Example 2 --- What an ALU actually is}

\begin{quote}
\textbf{An ALU is every function computed at once, and a multiplexer that picks
one.}
\end{quote}

All the units run in parallel, all the time, on the same inputs. The function code
does not tell the hardware what to compute --- it tells the \textbf{multiplexer}
which already-computed answer to forward. There is no ``choosing'' and then
``doing''; the doing already happened.

\textbf{So the only new structure in this lab is one you built in Lab 3.}

\subsection*{2.1 Two units and a mux, written structurally}

\begin{lstlisting}
U_ADD : entity work.alu_add port map (Ain => Ain, Bin => Bin, Rout => add_result, ...);
U_AND : entity work.alu_and port map (Ain => Ain, Bin => Bin, Rout => and_result);

result <= add_result when Func = '0' else and_result;    -- the multiplexer
\end{lstlisting}

\subsection*{2.2 The \code{case} \emph{is} the multiplexer}

At sixteen operations, writing sixteen entities and a 16-to-1 mux by hand is
pointless typing, so the assignment uses a \code{case}. That changes the spelling,
not the hardware:

\begin{quote}
\textbf{Open the RTL Viewer after the assignment and the synthesiser draws a
multiplexer.} A \code{case} is how you spell a mux when there are too many inputs
to wire by hand.
\end{quote}

\subsection*{2.3 The flag trap this example is built to prevent}

The adder computes its own \code{Zero} and \code{Neg}, and the example
\textbf{deliberately ignores them}:

\begin{lstlisting}
signal add_z, add_n, add_c : std_logic;   -- computed, then not used
...
Zero <= '1' when result = (result'range => '0') else '0';   -- on the SELECTED result
\end{lstlisting}

\begin{quote}
\textbf{Flags must describe the selected result, not one unit's result.} Compute
them \emph{after} the multiplexer, never inside a \code{case} branch.
\end{quote}

Concretely: \code{AND} of \code{0x0F0F0F0F} and \code{0xF0F0F0F0} is zero, so
\code{Zero} must be 1 --- but the adder's sum is \code{0xFFFFFFFF}, so a flag
taken from the adder reads 0. Nothing about the arithmetic is wrong; only the flag
is. In Lab 8 that broken flag becomes a branch that goes the wrong way.

\vspace{4pt}\hr

\section*{Example 3 --- The three shifts, and the one everybody gets wrong}

\begin{center}
\begin{tabular}{lll}
\toprule
& vacated bits filled with & correct when the value is \\
\midrule
\code{SLL} (= \code{SLA}) & \code{'0'} on the right & always --- there is only one sensible answer \\
\code{SRL} & \code{'0'} on the left & \textbf{unsigned} \\
\code{SRA} & \textbf{copies of the sign bit} & \textbf{signed} \\
\bottomrule
\end{tabular}
\end{center}

\textbf{Left shifts have only one sensible answer}, which is why \code{SLA} and
\code{SLL} are the same operation and the reference ALU handles them in one
branch.

\subsection*{3.1 Why \code{SRA} exists}

Shifting right by one should divide by two. Take $-8 =$ \code{0xFFFFFFF8}:

\begin{center}
\begin{tabular}{lll}
\toprule
& result & as a number \\
\midrule
\code{SRL} by 1 & \code{0x7FFFFFFC} & $+2\,147\,483\,644$ \quad --- $-8/2$ is not two billion \\
\code{SRA} by 1 & \code{0xFFFFFFFC} & $-4$ \quad --- correct \\
\bottomrule
\end{tabular}
\end{center}

Filling with \code{'0'} destroys the sign; filling with the sign bit preserves it.
\textbf{It is Lab 4's sign extension, applied while shifting.}

\subsection*{3.2 The bug this example exists to prevent}

\begin{lstlisting}
shift_right( unsigned(X), n )   -- this is SRL
shift_right(   signed(X), n )   -- this is SRA
\end{lstlisting}

\begin{quote}
\textbf{The operator name is the same. The type decides the behaviour.}
\end{quote}

So writing \code{SRA} with \code{unsigned()} silently gives you \code{SRL}. It
compiles, it runs, and \textbf{it is wrong only for negative operands} --- which
most casual test vectors do not include. Test with a negative number or you will
not find it.

\subsection*{3.3 The shift amount}

\begin{lstlisting}
signal amount : natural range 0 to 31;
amount <= to_integer(unsigned(Shft));
\end{lstlisting}

Five bits hold $0 \ldots 31$, exactly the useful range for a 32-bit shift:
nothing wasted, nothing out of range. The conversion is Lab 4's chain again ---
you cannot shift by a vector, because the amount must be an integer.

\textbf{And note which operand is shifted:} in this ALU shifts act on \code{Ain},
with the amount arriving on the separate \code{Shft} port. Assuming the amount
comes from \code{Bin} passes a home-made testbench and breaks the CPU in Lab 12.

\vspace{4pt}\hr

\section*{Into the assignment}

The reference ALU is these three examples at full size:

\begin{lstlisting}
when ALU_FUNC_SLL | ALU_FUNC_SLA =>
    result <= std_logic_vector(shift_left(unsigned(Ain), amount));
when ALU_FUNC_SRL =>
    result <= std_logic_vector(shift_right(unsigned(Ain), amount));
when ALU_FUNC_SRA =>
    result <= std_logic_vector(shift_right(signed(Ain), amount));   -- signed!
when ALU_FUNC_SOLS =>
    if signed(Ain) < signed(Bin) then result <= x"00000001";
    else result <= x"00000000"; end if;
when others =>
    result <= (others => '0');            -- assigns, so no latch
\end{lstlisting}

Four things to check in your own version:

\begin{enumerate}[leftmargin=2em,itemsep=2pt]
\item \textbf{\code{SRA} uses \code{signed()}, \code{SRL} uses \code{unsigned()}.}
Same operator, different fill.
\item \textbf{\code{SOLS} and \code{SOLU} use the \code{numeric\_std} comparison},
which is overflow-correct. Doing this in raw gates is the hard part of an ALU:
the hardware equation is $\text{SOLS} = R_{31} \oplus V$ with $R = A - B$, and
getting $V$ right is exactly what the library saves you.
\item \textbf{Flags are computed after the \code{case}}, on \code{result}, never
inside a branch.
\item \textbf{\code{when others} assigns}, so the process stays combinational ---
Lab 3's latch trap at sixteen times the scale.
\end{enumerate}

\vspace{6pt}
\hr
{\small \textbf{Forward.} The ALU is the widest component you will build and it
sets the CPU's critical path --- in Lab 12 the whole machine's $F_{max}$ is close
to this one block's. \code{Zero} and \code{Neg} leave here and become every
conditional branch in Lab 8.}

\end{document}
